% !TeX root = ../../main.tex
% =====================================================================
%  Prólogo
% =====================================================================
\chapter*{Prólogo}
\addcontentsline{toc}{chapter}{Prólogo}
\markboth{Prólogo}{Prólogo}

En un equipo de Formula Student el problema nunca ha sido la falta de talento.
El problema es que cada uno o dos años se marcha quien sabía por qué se tomó
cada decisión, y el que llega detrás vuelve a empezar desde cero. A veces incluso
peor: repitiendo un error que ya se había cometido y resuelto tres temporadas
antes, sin que quedara escrito en ninguna parte.

Este libro existe para romper ese ciclo. No es un manual de teoría estricto --- para eso
están los libros que se citan en la bibliografía, y son mejores que este. Es el
registro de \emph{cómo diseñamos nuestro coche}: qué decidimos, con qué datos,
qué salió bien, qué salió mal y qué haríamos distinto. La teoría está aquí solo
en la medida necesaria para entender esas decisiones. Habrá conceptos en los que se pueda
profundizar más debido a que tiene un gran impacto en el diseñor y otros que se tratan de
manera superficial porque son o demasiado complejos para explicarlos aquí, o porque no son
tan relevantes para el diseño en el punto en el que está el equipo. En un futuro se espera
que se vaya mejorando el nivel de detalle y de introducir más resultados aclaratorios de
los ensayos que se realizan en el equipo. De momento y como primera versión, se priorizará
la claridad y utilidad de la información, y se irá mejorando con el tiempo.

De ahí las tres reglas con las que está escrito:

\begin{enumerate}
  \item \textbf{Todo capítulo termina en algo que se usa.} Una plantilla, una
        hoja de cálculo, un checklist, un apartado del informe de diseño. Si un
        capítulo no deja herramienta, sobra.
  \item \textbf{Ningún número sin origen.} Cada valor que aparezca debe poder
        rastrearse hasta un ensayo, un cálculo o una referencia. Los números sin
        procedencia son los que arruinan un design event.
  \item \textbf{Se escribe mientras se hace, no al final.} Un capítulo redactado
        seis meses después de la decisión ya ha perdido lo importante: las
        alternativas que se descartaron y por qué. (en la primera versión del libro, 
        esto no se ha cumplido del todo, pero se espera que en futuras versiones se 
        cumpla)
\end{enumerate}

Este libro nunca estará terminado, y eso es exactamente lo que se pretende. Cada
temporada debe salir una versión mejor que la anterior. Si encuentras algo mal,
desactualizado o simplemente incompleto, la respuesta correcta no es avisar a
alguien: es corregirlo y firmar el cambio.

\vspace{1.5em}
\noindent\hfill\textit{El equipo \FUSAL}
